% ============================================================
% conclusion.tex — Sección VII: Conclusiones y Trabajo Futuro
% Estado: Preliminar (investigación en curso)
% ============================================================

\section{Conclusiones y Trabajo Futuro}
\label{sec:conclusion}

\subsection{Conclusiones preliminares}

El presente trabajo ha presentado los fundamentos teóricos y el diseño conceptual de \sirca{}, una arquitectura semi-autónoma de descubrimiento científico biomédico orientada a la investigación sobre plantas medicinales peruanas. La propuesta integra cinco componentes fundamentales —pipeline RAG biomédico, memoria vectorial persistente, agente semi-autónomo con retrieval híbrido, adquisición científica automatizada y generación contextualizada con grounding documental— en una infraestructura diseñada para abordar la fragmentación del conocimiento biomédico y las limitaciones de los sistemas de recuperación convencionales.

El análisis del estado del arte ha evidenciado que ningún sistema existente integra simultáneamente estas cinco capacidades orientadas al dominio etnobotánico, lo cual fundamenta la relevancia de la propuesta. Los fundamentos teóricos desarrollados —incluyendo la formalización del pipeline RAG, la representación vectorial biomédica, el retrieval híbrido, el reranking contextual y el grounding documental— proporcionan una base sólida para la implementación del sistema.

\subsection{Trabajo futuro}

Las líneas de trabajo futuro incluyen: (i)~la implementación del prototipo funcional de \sirca{} utilizando el stack tecnológico propuesto; (ii)~la construcción del corpus biomédico mediante adquisición automatizada; (iii)~la evaluación experimental cuantitativa según el marco definido en la Sección~\ref{sec:experiments}; (iv)~la exploración de grafos de conocimiento biomédico como complemento a la memoria vectorial; (v)~la incorporación de mecanismos de retroalimentación del usuario; y (vi)~la validación con investigadores activos en farmacognosia y etnobotánica peruana.

La viabilidad de la propuesta se sustenta en la disponibilidad de herramientas de código abierto, modelos biomédicos pre-entrenados y APIs de acceso programático a repositorios científicos, lo cual permite su implementación con recursos computacionales accesibles para entornos universitarios latinoamericanos.
